
It therefore suffices to show that, for every @@M0@@, one has
@@M1@@
where, as above, @@M2@@, @@M3@@. Now one has
@@M4@@
for every @@M5@@, so that
@@M6@@
But (6.2.10) and (6.5.3) give
@@M7@@
for every @@M8@@, whence
@@M9@@
with
@@M10@@
Now @@M11@@ is a section of @@M12@@, everywhere equal to the coefficient of @@M13@@ in @@M14@@, and is therefore zero. The result follows.

\sgaStableTarget{sga6.e10.sec.6}\section*{6. Appendix: Calculation of Determinants of Locally Free Sheaves}

\sgaStableTarget{sga6.e10.sec.6.1}\subsection*{6.1.}
First consider an arbitrary ringed topos @@M15@@. For every positive section @@M16@@ of the constant sheaf @@M17@@ and every @@M18@@-module @@M19@@, it is clear how to define the @@M20@@-modules
